% --- Archon physics formalization source begin ---
% archon:physics
% archon:covers PhyXMiniProblems/problem_phyx_mini_0714.lean
% archon:source-report reports/phyx_mini/problem_phyx_mini_0714.source.json

\chapter{Physics problem phyx\_mini\_0714}
\label{ch:PhyXMiniProblems_problem_phyx_mini_0714}

\paragraph{Problem source.}
\#\# Physical scenario

A 1000\textbackslash{} \textbackslash{}mathrm\{kg\} rocket is fired straight away from the surface of the earth. The rocket need to escape from the gravitational pull of the earth and never return. Assume a nonrotating earth.

\#\# Question

What speed does the rocket need?

\#\# Answer choices

- A: A: 12200m/s

- B: B: 13200m/s

- C: C: 11200m/s

- D: D: 14200m/s

\#\# Figure caption (auxiliary; use the image as primary evidence)

The image shows a diagram with two parts labeled "Before" and "After."

- **Before:**
  - There is a circle representing "Earth" with an arrow pointing from the center to the edge labeled \textbackslash{}( r\_1 = R\_e \textbackslash{}).
  - A rocket is depicted beside the Earth, moving away from it with an arrow indicating the velocity \textbackslash{}( v\_1 \textbackslash{}).

- **After:**
  - The rocket is shown further away without any thrust and two equations: 
    - \textbackslash{}( r\_2 = \textbackslash{}infty \textbackslash{}) indicating infinite distance.
    - \textbackslash{}( v\_2 = 0 \textbackslash{}, \textbackslash{}text\{m/s\} \textbackslash{}) indicating zero velocity.

The diagram illustrates a scenario of a rocket leaving Earth's influence, reaching infinite distance with a final velocity of zero.

\#\# Dataset metadata

Work and Energy; Physical Model Grounding Reasoning, Multi-Formula Reasoning

\paragraph{Recorded answer/context.}
C: C: 11200m/s

\paragraph{Figure/image path.}
/root/proposal\_for\_physic/science-mango/phyx\_mini\_run/phyx\_data/test\_image/714.png

\paragraph{Formalization target.}
create a compiling Lean file with sorry bodies at `PhyXMiniProblems/problem\_phyx\_mini\_0714.lean`. The Lean declarations must preserve the physical quantities, dimensions or dimensional roles, figure labels, governing-law hypotheses, and final relation expressed by this problem.
Use Mathlib/Physlib names found through LeanExplore where available. If a physics API is missing, introduce faithful local abstractions rather than scalar placeholder aliases.

\begin{theorem}[Physics formalization target]
\label{thm:physics:phyx_mini_0714:target}
\lean{PhyXMiniProblems.ProblemPhyXMini0714.problem_phyx_mini_0714}
\uses{def:physics:phyx-mini-0714:phyxminiproblems-problemphyxmini0714-massinkilograms, def:physics:phyx-mini-0714:phyxminiproblems-problemphyxmini0714-lengthinmeters, def:physics:phyx-mini-0714:phyxminiproblems-problemphyxmini0714-speedinmeterspersecond, def:physics:phyx-mini-0714:phyxminiproblems-problemphyxmini0714-gravitationalconstantinsi, def:physics:phyx-mini-0714:phyxminiproblems-problemphyxmini0714-surfaceescapesetup, def:physics:phyx-mini-0714:phyxminiproblems-problemphyxmini0714-matchesprimaryfigure, def:physics:phyx-mini-0714:phyxminiproblems-problemphyxmini0714-matchesproblemstatement, def:physics:phyx-mini-0714:phyxminiproblems-problemphyxmini0714-matchesstandardearthcalibration, def:physics:phyx-mini-0714:phyxminiproblems-problemphyxmini0714-hasphysicalescapeparameters, def:physics:phyx-mini-0714:phyxminiproblems-problemphyxmini0714-satisfiesnewtonianescapeenergylaws, def:physics:phyx-mini-0714:phyxminiproblems-problemphyxmini0714-recordedanswerchoice, def:physics:phyx-mini-0714:phyxminiproblems-problemphyxmini0714-isclosestanswerchoice}
The assigned autoformalize agent should translate this physics problem into Lean declarations in the covered file, with theorem and lemma proof bodies written as `by sorry`.
\end{theorem}
\begin{proof}
This is an autoformalization task, not a proof task. Produce statements that can later be proved by the physics prover without weakening the physical model.
\end{proof}

% --- Archon declaration topology begin ---
\section{Declaration topology}
\begin{definition}[gravitational Constant Dimension]
  \label{def:physics:phyx-mini-0714:phyxminiproblems-problemphyxmini0714-gravitationalconstantdimension}
  \lean{PhyXMiniProblems.ProblemPhyXMini0714.gravitationalConstantDimension}
  The physical dimension L³ M⁻¹ T⁻² of Newton's gravitational constant
\end{definition}
\begin{proof}
  This is the declared model definition of gravitational Constant Dimension.
\end{proof}
\begin{definition}[Mass Quantity]
  \label{def:physics:phyx-mini-0714:phyxminiproblems-problemphyxmini0714-massquantity}
  \lean{PhyXMiniProblems.ProblemPhyXMini0714.MassQuantity}
  A nonnegative, unit-independent physical mass
\end{definition}
\begin{proof}
  This is the declared model definition of Mass Quantity.
\end{proof}
\begin{definition}[Length Quantity]
  \label{def:physics:phyx-mini-0714:phyxminiproblems-problemphyxmini0714-lengthquantity}
  \lean{PhyXMiniProblems.ProblemPhyXMini0714.LengthQuantity}
  A nonnegative, unit-independent physical radial length
\end{definition}
\begin{proof}
  This is the declared model definition of Length Quantity.
\end{proof}
\begin{definition}[Speed Quantity]
  \label{def:physics:phyx-mini-0714:phyxminiproblems-problemphyxmini0714-speedquantity}
  \lean{PhyXMiniProblems.ProblemPhyXMini0714.SpeedQuantity}
  A nonnegative, unit-independent physical speed
\end{definition}
\begin{proof}
  This is the declared model definition of Speed Quantity.
\end{proof}
\begin{definition}[Angular Speed Magnitude Quantity]
  \label{def:physics:phyx-mini-0714:phyxminiproblems-problemphyxmini0714-angularspeedmagnitudequantity}
  \lean{PhyXMiniProblems.ProblemPhyXMini0714.AngularSpeedMagnitudeQuantity}
  A nonnegative angular-speed magnitude; radians are dimensionless
\end{definition}
\begin{proof}
  This is the declared model definition of Angular Speed Magnitude Quantity.
\end{proof}
\begin{definition}[Gravitational Constant Quantity]
  \label{def:physics:phyx-mini-0714:phyxminiproblems-problemphyxmini0714-gravitationalconstantquantity}
  \lean{PhyXMiniProblems.ProblemPhyXMini0714.GravitationalConstantQuantity}
  \uses{def:physics:phyx-mini-0714:phyxminiproblems-problemphyxmini0714-gravitationalconstantdimension}
  A unit-independent value of Newton's gravitational constant
\end{definition}
\begin{proof}
  This is the declared model definition of Gravitational Constant Quantity.
\end{proof}
\begin{definition}[Energy Quantity]
  \label{def:physics:phyx-mini-0714:phyxminiproblems-problemphyxmini0714-energyquantity}
  \lean{PhyXMiniProblems.ProblemPhyXMini0714.EnergyQuantity}
  A signed, unit-independent physical energy
\end{definition}
\begin{proof}
  This is the declared model definition of Energy Quantity.
\end{proof}
\begin{definition}[mass In Kilograms]
  \label{def:physics:phyx-mini-0714:phyxminiproblems-problemphyxmini0714-massinkilograms}
  \lean{PhyXMiniProblems.ProblemPhyXMini0714.massInKilograms}
  \uses{def:physics:phyx-mini-0714:phyxminiproblems-problemphyxmini0714-massquantity}
  Kilogram readout of a physical mass
\end{definition}
\begin{proof}
  This is the declared model definition of mass In Kilograms.
\end{proof}
\begin{definition}[length In Meters]
  \label{def:physics:phyx-mini-0714:phyxminiproblems-problemphyxmini0714-lengthinmeters}
  \lean{PhyXMiniProblems.ProblemPhyXMini0714.lengthInMeters}
  \uses{def:physics:phyx-mini-0714:phyxminiproblems-problemphyxmini0714-lengthquantity}
  Metre readout of a physical length
\end{definition}
\begin{proof}
  This is the declared model definition of length In Meters.
\end{proof}
\begin{definition}[speed In Meters Per Second]
  \label{def:physics:phyx-mini-0714:phyxminiproblems-problemphyxmini0714-speedinmeterspersecond}
  \lean{PhyXMiniProblems.ProblemPhyXMini0714.speedInMetersPerSecond}
  \uses{def:physics:phyx-mini-0714:phyxminiproblems-problemphyxmini0714-speedquantity}
  Metre-per-second readout of a physical speed
\end{definition}
\begin{proof}
  This is the declared model definition of speed In Meters Per Second.
\end{proof}
\begin{definition}[angular Speed In Radians Per Second]
  \label{def:physics:phyx-mini-0714:phyxminiproblems-problemphyxmini0714-angularspeedinradianspersecond}
  \lean{PhyXMiniProblems.ProblemPhyXMini0714.angularSpeedInRadiansPerSecond}
  \uses{def:physics:phyx-mini-0714:phyxminiproblems-problemphyxmini0714-angularspeedmagnitudequantity}
  Radian-per-second readout of an angular-speed magnitude
\end{definition}
\begin{proof}
  This is the declared model definition of angular Speed In Radians Per Second.
\end{proof}
\begin{definition}[gravitational Constant In SI]
  \label{def:physics:phyx-mini-0714:phyxminiproblems-problemphyxmini0714-gravitationalconstantinsi}
  \lean{PhyXMiniProblems.ProblemPhyXMini0714.gravitationalConstantInSI}
  \uses{def:physics:phyx-mini-0714:phyxminiproblems-problemphyxmini0714-gravitationalconstantquantity}
  SI readout of G, in m³ kg⁻¹ s⁻²
\end{definition}
\begin{proof}
  This is the declared model definition of gravitational Constant In SI.
\end{proof}
\begin{definition}[energy In Joules]
  \label{def:physics:phyx-mini-0714:phyxminiproblems-problemphyxmini0714-energyinjoules}
  \lean{PhyXMiniProblems.ProblemPhyXMini0714.energyInJoules}
  \uses{def:physics:phyx-mini-0714:phyxminiproblems-problemphyxmini0714-energyquantity}
  Joule readout of a signed physical energy
\end{definition}
\begin{proof}
  This is the declared model definition of energy In Joules.
\end{proof}
\begin{definition}[Earth Model]
  \label{def:physics:phyx-mini-0714:phyxminiproblems-problemphyxmini0714-earthmodel}
  \lean{PhyXMiniProblems.ProblemPhyXMini0714.EarthModel}
  \uses{def:physics:phyx-mini-0714:phyxminiproblems-problemphyxmini0714-massquantity, def:physics:phyx-mini-0714:phyxminiproblems-problemphyxmini0714-lengthquantity, def:physics:phyx-mini-0714:phyxminiproblems-problemphyxmini0714-angularspeedmagnitudequantity}
  The spherical Earth data relevant to the escape calculation
\end{definition}
\begin{proof}
  This is the declared model definition of Earth Model.
\end{proof}
\begin{definition}[Rocket Model]
  \label{def:physics:phyx-mini-0714:phyxminiproblems-problemphyxmini0714-rocketmodel}
  \lean{PhyXMiniProblems.ProblemPhyXMini0714.RocketModel}
  \uses{def:physics:phyx-mini-0714:phyxminiproblems-problemphyxmini0714-massquantity, def:physics:phyx-mini-0714:phyxminiproblems-problemphyxmini0714-speedquantity}
  The rocket's independent mass and before/after speed observables
\end{definition}
\begin{proof}
  This is the declared model definition of Rocket Model.
\end{proof}
\begin{definition}[Radial Location]
  \label{def:physics:phyx-mini-0714:phyxminiproblems-problemphyxmini0714-radiallocation}
  \lean{PhyXMiniProblems.ProblemPhyXMini0714.RadialLocation}
  \uses{def:physics:phyx-mini-0714:phyxminiproblems-problemphyxmini0714-lengthquantity}
  A radial location is either at a finite physical radius or at infinity
\end{definition}
\begin{proof}
  This is the declared model definition of Radial Location.
\end{proof}
\begin{definition}[Radial Direction]
  \label{def:physics:phyx-mini-0714:phyxminiproblems-problemphyxmini0714-radialdirection}
  \lean{PhyXMiniProblems.ProblemPhyXMini0714.RadialDirection}
  Radial direction of the rocket's initial velocity
\end{definition}
\begin{proof}
  This is the declared model definition of Radial Direction.
\end{proof}
\begin{definition}[Post Launch Propulsion]
  \label{def:physics:phyx-mini-0714:phyxminiproblems-problemphyxmini0714-postlaunchpropulsion}
  \lean{PhyXMiniProblems.ProblemPhyXMini0714.PostLaunchPropulsion}
  Propulsion state after the initial firing
\end{definition}
\begin{proof}
  This is the declared model definition of Post Launch Propulsion.
\end{proof}
\begin{definition}[Escape Phase]
  \label{def:physics:phyx-mini-0714:phyxminiproblems-problemphyxmini0714-escapephase}
  \lean{PhyXMiniProblems.ProblemPhyXMini0714.EscapePhase}
  The two panels explicitly named in the supplied image
\end{definition}
\begin{proof}
  This is the declared model definition of Escape Phase.
\end{proof}
\begin{definition}[Figure Radius Label]
  \label{def:physics:phyx-mini-0714:phyxminiproblems-problemphyxmini0714-figureradiuslabel}
  \lean{PhyXMiniProblems.ProblemPhyXMini0714.FigureRadiusLabel}
  Radius labels transcribed from the two panels
\end{definition}
\begin{proof}
  This is the declared model definition of Figure Radius Label.
\end{proof}
\begin{definition}[Figure Speed Label]
  \label{def:physics:phyx-mini-0714:phyxminiproblems-problemphyxmini0714-figurespeedlabel}
  \lean{PhyXMiniProblems.ProblemPhyXMini0714.FigureSpeedLabel}
  Speed labels transcribed from the two panels
\end{definition}
\begin{proof}
  This is the declared model definition of Figure Speed Label.
\end{proof}
\begin{definition}[Escape Figure]
  \label{def:physics:phyx-mini-0714:phyxminiproblems-problemphyxmini0714-escapefigure}
  \lean{PhyXMiniProblems.ProblemPhyXMini0714.EscapeFigure}
  \uses{def:physics:phyx-mini-0714:phyxminiproblems-problemphyxmini0714-escapephase, def:physics:phyx-mini-0714:phyxminiproblems-problemphyxmini0714-figureradiuslabel, def:physics:phyx-mini-0714:phyxminiproblems-problemphyxmini0714-figurespeedlabel}
  Literal qualitative and symbolic content shown in image 714.png
\end{definition}
\begin{proof}
  This is the declared model definition of Escape Figure.
\end{proof}
\begin{definition}[Surface Escape Setup]
  \label{def:physics:phyx-mini-0714:phyxminiproblems-problemphyxmini0714-surfaceescapesetup}
  \lean{PhyXMiniProblems.ProblemPhyXMini0714.SurfaceEscapeSetup}
  \uses{def:physics:phyx-mini-0714:phyxminiproblems-problemphyxmini0714-gravitationalconstantquantity, def:physics:phyx-mini-0714:phyxminiproblems-problemphyxmini0714-energyquantity, def:physics:phyx-mini-0714:phyxminiproblems-problemphyxmini0714-earthmodel, def:physics:phyx-mini-0714:phyxminiproblems-problemphyxmini0714-rocketmodel, def:physics:phyx-mini-0714:phyxminiproblems-problemphyxmini0714-radiallocation, def:physics:phyx-mini-0714:phyxminiproblems-problemphyxmini0714-radialdirection, def:physics:phyx-mini-0714:phyxminiproblems-problemphyxmini0714-postlaunchpropulsion, def:physics:phyx-mini-0714:phyxminiproblems-problemphyxmini0714-escapefigure}
  The independent physical observables for the marginal escape experiment. None of the speed fields is defined from an answer choice or from the desired escape-speed formula
\end{definition}
\begin{proof}
  This is the declared model definition of Surface Escape Setup.
\end{proof}
\begin{definition}[Matches Primary Figure]
  \label{def:physics:phyx-mini-0714:phyxminiproblems-problemphyxmini0714-matchesprimaryfigure}
  \lean{PhyXMiniProblems.ProblemPhyXMini0714.MatchesPrimaryFigure}
  \uses{def:physics:phyx-mini-0714:phyxminiproblems-problemphyxmini0714-speedinmeterspersecond, def:physics:phyx-mini-0714:phyxminiproblems-problemphyxmini0714-surfaceescapesetup}
  Primary-image transcription. The final zero-speed equality is a direct figure readout for the marginal-escape state, not the requested launch speed
\end{definition}
\begin{proof}
  This is the declared model definition of Matches Primary Figure.
\end{proof}
\begin{definition}[Matches Problem Statement]
  \label{def:physics:phyx-mini-0714:phyxminiproblems-problemphyxmini0714-matchesproblemstatement}
  \lean{PhyXMiniProblems.ProblemPhyXMini0714.MatchesProblemStatement}
  \uses{def:physics:phyx-mini-0714:phyxminiproblems-problemphyxmini0714-massinkilograms, def:physics:phyx-mini-0714:phyxminiproblems-problemphyxmini0714-angularspeedinradianspersecond, def:physics:phyx-mini-0714:phyxminiproblems-problemphyxmini0714-surfaceescapesetup}
  Numerical and qualitative facts stated in the problem prose
\end{definition}
\begin{proof}
  This is the declared model definition of Matches Problem Statement.
\end{proof}
\begin{definition}[Matches Standard Earth Calibration]
  \label{def:physics:phyx-mini-0714:phyxminiproblems-problemphyxmini0714-matchesstandardearthcalibration}
  \lean{PhyXMiniProblems.ProblemPhyXMini0714.MatchesStandardEarthCalibration}
  \uses{def:physics:phyx-mini-0714:phyxminiproblems-problemphyxmini0714-massinkilograms, def:physics:phyx-mini-0714:phyxminiproblems-problemphyxmini0714-lengthinmeters, def:physics:phyx-mini-0714:phyxminiproblems-problemphyxmini0714-gravitationalconstantinsi, def:physics:phyx-mini-0714:phyxminiproblems-problemphyxmini0714-surfaceescapesetup}
  Standard reference values needed to evaluate the recorded numerical choice. These data do not mention the launch speed or any answer choice
\end{definition}
\begin{proof}
  This is the declared model definition of Matches Standard Earth Calibration.
\end{proof}
\begin{definition}[Has Physical Escape Parameters]
  \label{def:physics:phyx-mini-0714:phyxminiproblems-problemphyxmini0714-hasphysicalescapeparameters}
  \lean{PhyXMiniProblems.ProblemPhyXMini0714.HasPhysicalEscapeParameters}
  \uses{def:physics:phyx-mini-0714:phyxminiproblems-problemphyxmini0714-massinkilograms, def:physics:phyx-mini-0714:phyxminiproblems-problemphyxmini0714-lengthinmeters, def:physics:phyx-mini-0714:phyxminiproblems-problemphyxmini0714-speedinmeterspersecond, def:physics:phyx-mini-0714:phyxminiproblems-problemphyxmini0714-gravitationalconstantinsi, def:physics:phyx-mini-0714:phyxminiproblems-problemphyxmini0714-surfaceescapesetup}
  Positivity conditions for the physical Newtonian branch
\end{definition}
\begin{proof}
  This is the declared model definition of Has Physical Escape Parameters.
\end{proof}
\begin{definition}[Satisfies Newtonian Escape Energy Laws]
  \label{def:physics:phyx-mini-0714:phyxminiproblems-problemphyxmini0714-satisfiesnewtonianescapeenergylaws}
  \lean{PhyXMiniProblems.ProblemPhyXMini0714.SatisfiesNewtonianEscapeEnergyLaws}
  \uses{def:physics:phyx-mini-0714:phyxminiproblems-problemphyxmini0714-massinkilograms, def:physics:phyx-mini-0714:phyxminiproblems-problemphyxmini0714-lengthinmeters, def:physics:phyx-mini-0714:phyxminiproblems-problemphyxmini0714-speedinmeterspersecond, def:physics:phyx-mini-0714:phyxminiproblems-problemphyxmini0714-gravitationalconstantinsi, def:physics:phyx-mini-0714:phyxminiproblems-problemphyxmini0714-energyinjoules, def:physics:phyx-mini-0714:phyxminiproblems-problemphyxmini0714-surfaceescapesetup}
  Newtonian work-energy model for unpowered radial flight. The kinetic-energy relations are K = mv²/2; the finite-radius potential is U = -GMm/r; the potential tends to zero at infinity; and total mechanical energy is conserved while coasting. No field states v₁ = sqrt (2GM/R) or selects a displayed answer
\end{definition}
\begin{proof}
  This is the declared model definition of Satisfies Newtonian Escape Energy Laws.
\end{proof}
\begin{definition}[Answer Choice]
  \label{def:physics:phyx-mini-0714:phyxminiproblems-problemphyxmini0714-answerchoice}
  \lean{PhyXMiniProblems.ProblemPhyXMini0714.AnswerChoice}
  Labels of the four speed choices displayed with the problem
\end{definition}
\begin{proof}
  This is the declared model definition of Answer Choice.
\end{proof}
\begin{definition}[speed In Meters Per Second]
  \label{def:physics:phyx-mini-0714:phyxminiproblems-problemphyxmini0714-answerchoice-speedinmeterspersecond}
  \lean{PhyXMiniProblems.ProblemPhyXMini0714.AnswerChoice.speedInMetersPerSecond}
  \uses{def:physics:phyx-mini-0714:phyxminiproblems-problemphyxmini0714-speedinmeterspersecond, def:physics:phyx-mini-0714:phyxminiproblems-problemphyxmini0714-answerchoice}
  Metre-per-second value printed beside each answer label
\end{definition}
\begin{proof}
  This is the declared model definition of speed In Meters Per Second.
\end{proof}
\begin{definition}[recorded Answer Choice]
  \label{def:physics:phyx-mini-0714:phyxminiproblems-problemphyxmini0714-recordedanswerchoice}
  \lean{PhyXMiniProblems.ProblemPhyXMini0714.recordedAnswerChoice}
  \uses{def:physics:phyx-mini-0714:phyxminiproblems-problemphyxmini0714-answerchoice}
  The answer label recorded in the supplied dataset
\end{definition}
\begin{proof}
  This is the declared model definition of recorded Answer Choice.
\end{proof}
\begin{definition}[Is Closest Answer Choice]
  \label{def:physics:phyx-mini-0714:phyxminiproblems-problemphyxmini0714-isclosestanswerchoice}
  \lean{PhyXMiniProblems.ProblemPhyXMini0714.IsClosestAnswerChoice}
  \uses{def:physics:phyx-mini-0714:phyxminiproblems-problemphyxmini0714-speedinmeterspersecond, def:physics:phyx-mini-0714:phyxminiproblems-problemphyxmini0714-answerchoice}
  A choice is closest when its displayed speed has no greater absolute error than any other displayed speed. This avoids falsely identifying the rounded 11200 m/s choice with the exact speed from the calibrated constants
\end{definition}
\begin{proof}
  This is the declared model definition of Is Closest Answer Choice.
\end{proof}
% --- Archon declaration topology end ---
% --- Archon physics formalization source end ---
